
% 150-200 Palabras
% Hacer hasta el Final
        Usando mediciones de voltaje en el cuero cabelludo (EEG por sus siglas en ingles) o mediciones de campo magnético en la superficie de la cabeza (MEG por sus siglas en ingles) se busca estimar su fuente de emisión. Esto implica trabajar hacia atrás del sensor (Problema Inverso) para determinar las fuentes intercraneales esto nos proveerá con características espacio-temporales de la fuente neuronal.

    \subsection{Métodos Actuales}
    Los métodos actuales de localización de fuentes cerebrales a partir de señales de EEG o MEG se dividen principalmente en dos grandes categorías: los modelos paramétricos, que suponen un número reducido y definido de fuentes activas (como el dipole fitting y MUSIC), y los modelos distribuidos, que estiman una distribución continua de corriente en el volumen cerebral (como MNE, LORETA, sLORETA, VARETA y LAURA). Estos enfoques han permitido importantes avances en el análisis funcional del cerebro, especialmente en aplicaciones clínicas como la epilepsia o la neurorehabilitación.

    Sin embargo, todos estos métodos comparten una dificultad inherente: el problema inverso es fundamentalmente mal condicionado. Es decir, existen múltiples configuraciones internas del cerebro que pueden generar las mismas señales observadas en la superficie. Esta ambigüedad limita la precisión espacial y temporal de las estimaciones y requiere incorporar supuestos o restricciones adicionales, lo cual puede introducir sesgos o reducir la generalización del modelo. Resolver esta problemática sigue siendo uno de los principales desafíos en el campo de la neuroimagen no invasiva.